\documentclass[11pt,a4paper]{article}

% ---------- Encodage et langue ----------
\usepackage[utf8]{inputenc}
\usepackage[T1]{fontenc}
% \usepackage{lmodern}  % décommenter si disponible (police vectorielle)
\usepackage[french]{babel}
\usepackage[a4paper,margin=2.3cm]{geometry}

% ---------- Mathématiques ----------
\usepackage{amsmath,amssymb,amsthm}
\usepackage{mathtools}
\usepackage{enumitem}
\usepackage[most]{tcolorbox}
\usepackage[colorlinks=true,linkcolor=blue!60!black,urlcolor=blue!60!black]{hyperref}

% ---------- Macros ----------
\newcommand{\R}{\mathbb{R}}
\newcommand{\N}{\mathbb{N}}
\newcommand{\Z}{\mathbb{Z}}
\newcommand{\Q}{\mathbb{Q}}
\newcommand{\C}{\mathbb{C}}
\newcommand{\K}{\mathbb{K}}
\newcommand{\Sun}{\mathbb{S}^1}
\newcommand{\interieur}[1]{\mathring{#1}}
\newcommand{\adh}[1]{\overline{#1}}
\newcommand{\Int}{\operatorname{int}}
\newcommand{\diam}{\operatorname{diam}}
\newcommand{\dist}{\operatorname{dist}}
\newcommand{\Vect}{\operatorname{Vect}}
\newcommand{\Bf}{B_f}
\newcommand{\norm}[1]{\left\|#1\right\|}
\newcommand{\abs}[1]{\left|#1\right|}
\newcommand{\opnorm}[1]{|\!|\!|#1|\!|\!|}
\newcommand{\ps}[2]{\left\langle #1,#2\right\rangle}
\newcommand{\Lc}{\mathcal{L}}
\newcommand{\Cz}{\mathcal{C}^0}
\newcommand{\Cinf}{\mathcal{C}^\infty}
\newcommand{\Cc}{\mathcal{C}^0_c}
\newcommand{\ie}{c'est-à-dire}
\DeclareMathOperator{\Ker}{Ker}
\DeclareMathOperator{\Ima}{Im}
\DeclareMathOperator{\supp}{supp}
\DeclareMathOperator{\sgn}{sgn}
\renewcommand{\Re}{\operatorname{Re}}

% ---------- Environnements ----------
\newtcolorbox{exo}[1]{
  enhanced, breakable,
  colback=black!3, colframe=black!70,
  fonttitle=\bfseries, title={#1},
  boxrule=0.6pt, arc=1pt, left=6pt, right=6pt
}
\newtcolorbox{pointcle}{
  enhanced, breakable,
  colback=blue!4, colframe=blue!55!black,
  fonttitle=\bfseries, title={Point clé},
  boxrule=0.8pt, arc=2pt, left=6pt, right=6pt
}
\newtcolorbox{piege}{
  enhanced, breakable,
  colback=red!4, colframe=red!60!black,
  fonttitle=\bfseries, title={Piège / difficulté},
  boxrule=0.8pt, arc=2pt, left=6pt, right=6pt
}
\newtcolorbox{commentaire}{
  enhanced, breakable,
  colback=gray!8, colframe=gray!55!black,
  fonttitle=\bfseries, title={Commentaire},
  boxrule=0.6pt, arc=2pt, left=6pt, right=6pt
}
\newtcolorbox{methode}{
  enhanced, breakable,
  colback=green!4, colframe=green!40!black,
  fonttitle=\bfseries, title={Méthode},
  boxrule=0.8pt, arc=2pt, left=6pt, right=6pt
}

\theoremstyle{definition}
\newtheorem*{solution}{Solution}

\title{\textbf{4MA305 -- Bases d'analyse fonctionnelle}\\[2mm]
\Large Révision des chapitres 1 à 3 -- Corrigé détaillé}
\author{}
\date{Année universitaire 2026--2027}

\begin{document}
\maketitle

\begin{commentaire}
Feuille de révision des chapitres~1 à~3 (espaces métriques, espaces vectoriels normés, espaces de Banach), composée dans l'optique de l'examen sur le modèle des sections « Vers l'examen » des TD4, TD5 et TD6. Elle ne figure pas parmi les feuilles officielles. Les Exercices~1 et~2 font restituer les preuves du polycopié le plus souvent demandées en question de cours ; les Exercices~3 et~4 sont des vrai ou faux sur les pièges de ces trois chapitres. La numérotation des énoncés cités est celle du polycopié ; les encadrés « Méthode » et « Piège » résument ce qu'il faut retenir.
\end{commentaire}

\section*{Questions de cours}

% =====================================================================
\begin{exo}{Exercice 1 -- Questions de cours : chapitre 1}
\emph{À traiter sans document : énoncé complet, puis démonstration.}
\begin{enumerate}[label=(\alph*)]
  \item Énoncer et démontrer le théorème de prolongement des applications uniformément continues (Théorème~1.1).
  \item Énoncer et démontrer le lemme de Baire (Lemme~1.1).
  \item Énoncer et démontrer le théorème de Heine (Théorème~1.3).
  \item Montrer que dans un espace métrique complet, une partie est compacte si et seulement si elle est fermée et précompacte (Corollaire~1.4 ; on redémontrera les Propositions~1.13 et 1.14).
  \item Énoncer et démontrer le théorème de point fixe de Picard (Exercice~1.16).
\end{enumerate}
\end{exo}

\begin{solution}
\noindent\textbf{(a)} \emph{Énoncé.} Soient $(X,d_X)$ un espace métrique, $(Y,d_Y)$ un espace métrique \emph{complet}, $D$ une partie dense de $X$ et $f:D\to Y$ uniformément continue. Il existe une unique application continue $\tilde f:X\to Y$ prolongeant $f$ ; elle est uniformément continue.

\emph{Unicité.} Deux applications continues qui coïncident sur une partie dense coïncident partout (caractérisation séquentielle).

\emph{Construction.} Soit $x\in X$ et $(d_n)$ une suite de $D$ convergeant vers $x$ (densité). Fixons $\varepsilon>0$ et $\delta>0$ donné par l'uniforme continuité : $d_X(a,b)<\delta\Rightarrow d_Y(f(a),f(b))<\varepsilon$ pour $a,b\in D$. La suite $(d_n)$ est de Cauchy : il existe $N$ avec $d_X(d_n,d_m)<\delta$ pour $n,m\ge N$, donc $d_Y(f(d_n),f(d_m))<\varepsilon$. Ainsi $(f(d_n))$ est de Cauchy dans $Y$ complet : elle converge. Si $(e_n)$ est une autre suite de $D$ tendant vers $x$, la suite alternée $d_0,e_0,d_1,e_1,\dots$ tend aussi vers $x$, donc son image converge, et les deux suites extraites $(f(d_n))$ et $(f(e_n))$ ont la même limite. On pose $\tilde f(x)$ cette limite commune ; pour $x\in D$ la suite constante donne $\tilde f(x)=f(x)$.

\emph{Uniforme continuité.} Soient $x_1,x_2\in X$ avec $d_X(x_1,x_2)<\delta/3$, et $(a_n)$, $(b_n)$ dans $D$ tendant vers $x_1$, $x_2$. Pour $n$ assez grand, $d_X(a_n,b_n)\le d_X(a_n,x_1)+d_X(x_1,x_2)+d_X(x_2,b_n)<\delta$, donc $d_Y(f(a_n),f(b_n))<\varepsilon$ ; en passant à la limite (continuité de $d_Y$), $d_Y(\tilde f(x_1),\tilde f(x_2))\le\varepsilon$.

\medskip
\noindent\textbf{(b)} \emph{Énoncé.} Dans un espace métrique complet : (i) toute intersection dénombrable d'ouverts denses est dense ; (ii) toute réunion dénombrable de fermés d'intérieur vide est d'intérieur vide.

\emph{Preuve.} (i) et (ii) sont équivalents par passage au complémentaire (un fermé est d'intérieur vide si et seulement si son complémentaire est un ouvert dense). Montrons (i). Soient $(O_n)$ des ouverts denses, $O:=\bigcap_nO_n$ et $U$ un ouvert non vide ; il s'agit de montrer $U\cap O\ne\emptyset$ (Proposition~1.4). L'ouvert $U\cap O_0$ est non vide ($O_0$ dense) : il contient une boule fermée $\Bf(a_0,r_0)$ avec $0<r_0\le1$. L'ouvert $B(a_0,r_0)\cap O_1$ est non vide : il contient $\Bf(a_1,r_1)$ avec $0<r_1\le\frac12$. Par récurrence on construit $a_n$ et $0<r_n\le2^{-n}$ avec
\[
\Bf(a_n,r_n)\subset B(a_{n-1},r_{n-1})\cap O_n .
\]
Les fermés $F_n:=\Bf(a_n,r_n)$ sont emboîtés. Pour $m\ge n$, $a_m\in F_n$, donc $d(a_n,a_m)\le2^{-n}$ : la suite $(a_n)$ est de Cauchy, elle converge vers un $a$ (complétude). Pour tout $n$, la suite $(a_m)_{m\ge n}$ est dans le fermé $F_n$, donc $a\in F_n$. Ainsi $a\in F_0\subset U$ et $a\in F_n\subset O_n$ pour tout $n$ : $a\in U\cap O$.

\medskip
\noindent\textbf{(c)} \emph{Énoncé.} Soient $X$ un espace métrique compact, $Y$ un espace métrique et $f:X\to Y$ continue. Alors $f(X)$ est compact et $f$ est uniformément continue.

\emph{Compacité de l'image.} Soit $(y_n)$ dans $f(X)$, $y_n=f(x_n)$. Par compacité, $x_{\varphi(n)}\to x\in X$, et par continuité $y_{\varphi(n)}\to f(x)\in f(X)$.

\emph{Uniforme continuité.} Par l'absurde : il existerait $\varepsilon>0$ et, pour tout $n\ge1$, des points $a_n,b_n$ avec $d_X(a_n,b_n)\le\frac1n$ et $d_Y(f(a_n),f(b_n))\ge\varepsilon$. Par compacité (deux extractions successives), $a_{\varphi(n)}\to a$ et $b_{\varphi(n)}\to b$ ; comme $d_X(a_n,b_n)\to0$, $a=b$. Par continuité de $f$ et de $d_Y$, $d_Y(f(a_{\varphi(n)}),f(b_{\varphi(n)}))\to d_Y(f(a),f(a))=0$, ce qui contredit $\ge\varepsilon$.

\medskip
\noindent\textbf{(d)} Soit $A$ une partie d'un espace métrique complet $(X,d)$.

\emph{Compacte $\Rightarrow$ fermée et précompacte.} \emph{Fermée :} si $a_n\in A$ et $a_n\to x$, une sous-suite converge dans $A$, nécessairement vers $x$ : $x\in A$. \emph{Précompacte (Proposition~1.13) :} par contraposée, si pour un $\varepsilon>0$ aucune famille finie de boules de rayon $\varepsilon$ ne recouvre $A$, on construit par récurrence $x_n\in A\setminus\bigcup_{k<n}B(x_k,\varepsilon)$ ; alors $d(x_n,x_m)\ge\varepsilon$ pour $n\ne m$, aucune sous-suite n'est de Cauchy, donc aucune ne converge : $A$ n'est pas compacte.

\emph{Fermée et précompacte $\Rightarrow$ compacte.} Soit $(x_n)$ dans $A$. \emph{Extraction d'une sous-suite de Cauchy (Proposition~1.14) :} $A$ est recouverte par un nombre fini de boules de rayon $1$ ; l'une d'elles, $B_0$, contient $x_n$ pour une infinité d'indices. Ces termes sont recouverts par un nombre fini de boules de rayon $\frac12$ ; l'une d'elles, $B_1$, en contient une infinité. On construit ainsi des boules $B_k$ de rayon $2^{-k}$ telles que $B_0\cap\dots\cap B_k$ contienne $x_n$ pour une infinité de $n$, puis une extractrice $\varphi$ avec $x_{\varphi(k)}\in B_0\cap\dots\cap B_k$. Pour $j\ge k$, $x_{\varphi(j)}$ et $x_{\varphi(k)}$ sont dans $B_k$, donc $d(x_{\varphi(j)},x_{\varphi(k)})\le2^{1-k}$ : la suite extraite est de Cauchy. \emph{Conclusion :} $X$ étant complet, elle converge, et sa limite est dans $A$ fermé.

\medskip
\noindent\textbf{(e)} \emph{Énoncé.} Soit $(X,d)$ un espace métrique complet non vide et $f:X\to X$ une application $k$-lipschitzienne avec $k<1$. Alors $f$ a un unique point fixe.

\emph{Unicité.} Si $f(x)=x$ et $f(y)=y$, alors $d(x,y)=d(f(x),f(y))\le k\,d(x,y)$, donc $(1-k)d(x,y)\le0$ et $x=y$.

\emph{Existence.} Soit $x_0\in X$ et $x_{n+1}:=f(x_n)$. Par récurrence, $d(x_{n+1},x_n)\le k^nd(x_1,x_0)$, donc pour $m>n$,
\[
d(x_m,x_n)\le\sum_{j=n}^{m-1}d(x_{j+1},x_j)\le\frac{k^n}{1-k}\,d(x_1,x_0)\xrightarrow[n\to\infty]{}0 .
\]
La suite est de Cauchy, donc converge vers un $x\in X$. Comme $f$ est continue (lipschitzienne), $f(x)=\lim f(x_n)=\lim x_{n+1}=x$.
\end{solution}

\begin{methode}
\textbf{Le rôle de la complétude, à chaque fois.} Dans (a), (b), (d) et (e), la preuve suit le même schéma : fabriquer une suite de Cauchy (par uniforme continuité, par boules emboîtées de rayons $2^{-n}$, par précompacité, par contraction), la faire converger grâce à la complétude, puis identifier la limite. À l'examen, dire explicitement \emph{où} la complétude intervient : c'est l'hypothèse que l'on vous demandera de justifier. Dans (b), ne pas oublier de prendre des boules \emph{fermées} incluses dans la boule \emph{ouverte} précédente : c'est ce qui garantit que la limite reste dans tous les $F_n$.
\end{methode}

% =====================================================================
\begin{exo}{Exercice 2 -- Questions de cours : chapitres 2 et 3}
\emph{À traiter sans document : énoncé complet, puis démonstration.}
\begin{enumerate}[label=(\alph*)]
  \item Énoncer et démontrer la caractérisation des applications linéaires continues entre EVN (Théorème~2.4).
  \item Énoncer et démontrer l'équivalence des normes sur $\K^d$ (Théorème~2.1).
  \item Montrer qu'un EVN est de Banach si et seulement si toute série absolument convergente y converge (Proposition~3.1 et Exercice~3.1).
  \item Énoncer et démontrer le théorème de Riesz (Théorème~3.1), en passant par le Lemme~3.1.
  \item Montrer que $\Lc(E,F)$ est un espace de Banach dès que $F$ en est un (Proposition~3.4).
  \item Énoncer et démontrer le théorème de Banach--Steinhaus (Théorème~3.7).
  \item Énoncer le théorème de l'application ouverte (Théorème~3.8, admis ici). En déduire le Corollaire~3.4 (continuité automatique de l'inverse), puis énoncer et démontrer le théorème du graphe fermé (Remarque~3.8).
\end{enumerate}
\end{exo}

\begin{solution}
\noindent\textbf{(a)} \emph{Énoncé.} Soient $E,F$ deux EVN et $T:E\to F$ linéaire. Sont équivalentes : (i) $T$ est lipschitzienne ; (ii) $T$ est continue en $0$ ; (iii) $T$ est bornée sur $B_E(0,1)$ ; (iv) il existe $C>0$ tel que $\norm{T(x)}_F\le C\norm x_E$ pour tout $x$.

\emph{Preuve.} (i)$\Rightarrow$(ii) : immédiat. (ii)$\Rightarrow$(iii) : il existe $\delta>0$ tel que $\norm x<\delta$ entraîne $\norm{T(x)}<1$ ; pour $\norm x<1$, $\norm{T(x)}=\delta^{-1}\norm{T(\delta x)}<\delta^{-1}$. (iii)$\Rightarrow$(iv) : soit $M$ une borne de $\norm{T}$ sur $B_E(0,1)$ ; pour $x\ne0$, $\frac x{2\norm x}\in B_E(0,1)$, donc $\norm{T(x)}=2\norm x\,\big\|T\big(\frac x{2\norm x}\big)\big\|\le2M\norm x$. (iv)$\Rightarrow$(i) : par linéarité, $\norm{T(x)-T(y)}=\norm{T(x-y)}\le C\norm{x-y}$.

\medskip
\noindent\textbf{(b)} \emph{Énoncé.} Sur $\K^d$, toutes les normes sont équivalentes.

\emph{Preuve.} Il suffit de comparer une norme $N$ à $\norm\cdot_\infty$. \emph{Majoration :} avec la base canonique $(e_k)$, $N(x)\le\sum_k|x_k|N(e_k)\le\beta\norm x_\infty$, où $\beta:=\sum_kN(e_k)$. En particulier $|N(x)-N(y)|\le N(x-y)\le\beta\norm{x-y}_\infty$ : $N$ est continue sur $(\K^d,\norm\cdot_\infty)$. \emph{Minoration :} la sphère $S_\infty:=\{\norm z_\infty=1\}$ est fermée et bornée dans $(\K^d,\norm\cdot_\infty)$, donc compacte (Corollaire~1.2). La fonction continue $N$ y atteint son minimum $\alpha$ (Corollaire~1.3), et $\alpha>0$ puisque $N$ ne s'annule pas sur $S_\infty$. Par homogénéité, pour $x\ne0$, $N(x)=\norm x_\infty N(x/\norm x_\infty)\ge\alpha\norm x_\infty$. Donc $\alpha\norm\cdot_\infty\le N\le\beta\norm\cdot_\infty$.

\medskip
\noindent\textbf{(c)} \emph{Sens direct (Proposition~3.1).} Soit $E$ de Banach et $\sum_k\norm{x_k}<\infty$. Pour $n<m$, $\norm{S_m-S_n}\le\sum_{k=n+1}^m\norm{x_k}\le\sum_{k>n}\norm{x_k}\to0$ : $(S_n)$ est de Cauchy, donc converge.

\emph{Réciproque (Exercice~3.1).} Supposons que toute série absolument convergente converge, et soit $(x_n)$ de Cauchy dans $E$. On construit une extractrice $\varphi$ telle que $\norm{x_p-x_q}\le2^{-k}$ pour tous $p,q\ge\varphi(k)$. Alors $\sum_k\norm{x_{\varphi(k+1)}-x_{\varphi(k)}}\le\sum_k2^{-k}<\infty$, donc la série $\sum_k(x_{\varphi(k+1)}-x_{\varphi(k)})$ converge ; ses sommes partielles valant $x_{\varphi(K+1)}-x_{\varphi(0)}$, la sous-suite $(x_{\varphi(k)})$ converge vers un $x$. Enfin une suite de Cauchy qui a une sous-suite convergente converge : $\norm{x_n-x}\le\norm{x_n-x_{\varphi(k)}}+\norm{x_{\varphi(k)}-x}$, petit pour $n$ et $k$ grands.

\medskip
\noindent\textbf{(d)} \emph{Énoncé.} Un EVN dont la boule unité fermée est compacte est de dimension finie.

\emph{Lemme~3.1.} S'il existe $0<\delta<1$ et $x_1,\dots,x_N\in E$ tels que la sphère unité soit contenue dans $\bigcup_kB(x_k,\delta)$, alors $(x_k)$ engendre $E$.

\emph{Preuve du lemme.} $F:=\Vect(x_1,\dots,x_N)$ est de dimension finie, donc fermé (Proposition~3.2). Supposons qu'il existe $z\in E\setminus F$ ; en particulier $z\ne0$. Il existe $k$ avec $\norm{z/\norm z-x_k}<\delta$, soit, pour $z_1:=\norm z\,x_k\in F$, $\norm{z-z_1}\le\delta\norm z$. Le vecteur $z-z_1$ est encore hors de $F$ : le même argument donne $z_2\in F$ avec $\norm{z-z_1-z_2}\le\delta\norm{z-z_1}\le\delta^2\norm z$. Par récurrence, il existe $z_n\in F$ avec $\norm{z-(z_1+\dots+z_n)}\le\delta^n\norm z\to0$. Donc $z$ est limite d'une suite de $F$ fermé : $z\in F$, contradiction.

\emph{Preuve du théorème.} La boule unité fermée, compacte, est précompacte (Proposition~1.13) : elle est recouverte par un nombre fini de boules $B(x_k,\frac12)$, et la sphère unité aussi. Le lemme avec $\delta=\frac12$ donne $E=\Vect(x_1,\dots,x_N)$.

\medskip
\noindent\textbf{(e)} Soit $(T_n)$ de Cauchy dans $\Lc(E,F)$, avec $F$ de Banach. Pour $x\in E$, $\norm{T_n(x)-T_m(x)}\le\opnorm{T_n-T_m}\norm x$ : la suite $(T_n(x))$ est de Cauchy dans $F$, donc converge ; on note $T(x)$ sa limite. \emph{Linéarité :} passage à la limite dans $T_n(\lambda x+y)=\lambda T_n(x)+T_n(y)$. \emph{Convergence en norme :} soit $\varepsilon>0$ et $N$ tel que $\opnorm{T_n-T_m}\le\varepsilon$ pour $n,m\ge N$. Pour $x\in E$ et $n\ge N$, en faisant $m\to\infty$ dans $\norm{T_n(x)-T_m(x)}\le\varepsilon\norm x$, on obtient $\norm{T_n(x)-T(x)}\le\varepsilon\norm x$. Donc $T_N-T$ est continue, $T=T_N-(T_N-T)\in\Lc(E,F)$, et $\opnorm{T_n-T}\le\varepsilon$ pour $n\ge N$.

\medskip
\noindent\textbf{(f)} \emph{Énoncé.} Soient $E$ un espace de \emph{Banach}, $F$ un EVN et $(T_i)_{i\in I}$ une famille de $\Lc(E,F)$. Si $\sup_i\norm{T_i(x)}_F<\infty$ pour tout $x\in E$, alors $\sup_i\opnorm{T_i}<\infty$.

\emph{Preuve.} Pour $n\in\N$, soit $K_n:=\{x\in E:\ \forall i,\ \norm{T_i(x)}\le n\}=\bigcap_i(\norm\cdot_F\circ T_i)^{-1}([0,n])$ : c'est un fermé (intersection d'images réciproques de fermés par des applications continues). L'hypothèse donne $E=\bigcup_nK_n$. Par le lemme de Baire dans $E$ complet, les $K_n$ ne peuvent pas être tous d'intérieur vide : il existe $N$, $a\in E$ et $r>0$ avec $\Bf(a,r)\subset K_N$. Pour $\norm z\le1$ et $i\in I$, $a+rz\in K_N$ et $a\in K_N$, donc
\[
\norm{T_i(z)}=\frac1r\norm{T_i(a+rz)-T_i(a)}\le\frac{2N}r .
\]
Ainsi $\opnorm{T_i}\le2N/r$ pour tout $i$.

\medskip
\noindent\textbf{(g)} \emph{Application ouverte.} Soient $E,F$ deux espaces de Banach et $T\in\Lc(E,F)$ surjective. Alors $T$ est ouverte : l'image de tout ouvert de $E$ est un ouvert de $F$.

\emph{Corollaire~3.4.} Si de plus $T$ est bijective, $T^{-1}$ est continue : pour tout ouvert $U$ de $E$,
\[
(T^{-1})^{-1}(U)=T(U)\quad\text{est ouvert,}
\]
ce qui est la caractérisation topologique de la continuité (Proposition~1.5).

\emph{Graphe fermé.} Soient $E,F$ deux espaces de Banach et $T:E\to F$ linéaire. Alors $T$ est continue si et seulement si son graphe $G(T):=\{(x,T(x)):x\in E\}$ est fermé dans $E\times F$.

\emph{Preuve.} ($\Rightarrow$) Si $(x_n,T(x_n))\to(x,y)$, alors $T(x_n)\to T(x)$ par continuité, donc $y=T(x)$. ($\Leftarrow$) Posons $N(x):=\norm x_E+\norm{T(x)}_F$ : c'est une norme sur $E$. \emph{$(E,N)$ est complet :} si $(x_n)$ est de Cauchy pour $N$, alors $(x_n)$ est de Cauchy dans $E$ et $(T(x_n))$ dans $F$ ; elles convergent, vers $x$ et $y$. Le graphe étant fermé, $y=T(x)$, donc $N(x_n-x)=\norm{x_n-x}_E+\norm{T(x_n)-y}_F\to0$. \emph{Conclusion :} l'identité $(E,N)\to(E,\norm\cdot_E)$ est linéaire, bijective et continue ($\norm x_E\le N(x)$) entre deux espaces de Banach ; par le Corollaire~3.4 son inverse est continu : il existe $C$ avec $N(x)\le C\norm x_E$, d'où $\norm{T(x)}_F\le C\norm x_E$.
\end{solution}

\begin{methode}
\textbf{Trois théorèmes, un seul moteur.} Banach--Steinhaus, l'application ouverte et le graphe fermé sont trois conséquences du lemme de Baire, et les trois demandent la complétude de l'espace de \emph{départ} (les deux derniers aussi celle de l'espace d'arrivée). Dans (f), le plan est : des fermés $K_n$ qui recouvrent $E$, Baire pour en trouver un d'intérieur non vide, puis translation et homothétie pour revenir à la boule unité. Dans (g), l'idée à retenir est la \emph{norme du graphe} $N(x)=\norm x+\norm{T(x)}$ : le graphe fermé dit exactement qu'elle rend $E$ complet, et deux normes comparables qui rendent un espace complet sont équivalentes (Exemple~3.6). En pratique, le graphe fermé sert à montrer la continuité en ne vérifiant que « si $x_n\to x$ \emph{et} $T(x_n)\to y$, alors $y=T(x)$ », ce qui est bien plus faible que la continuité.
\end{methode}

\section*{Vrai ou faux}

% =====================================================================
\begin{exo}{Exercice 3 -- Vrai ou faux ? Chapitre 1}
Pour chaque affirmation, dire si elle est vraie ou fausse, et justifier par une preuve courte ou un contre-exemple. $(X,d)$ et $(Y,d')$ désignent des espaces métriques.
\begin{enumerate}[label=\arabic*.]
  \item L'adhérence de la boule ouverte $B(a,r)$ est la boule fermée $\Bf(a,r)$.
  \item Dans un espace métrique complet, toute partie fermée et bornée est compacte.
  \item L'image réciproque d'un compact par une application continue est compacte.
  \item L'image d'un fermé par une application continue est fermée.
  \item Une application continue transforme les suites de Cauchy en suites de Cauchy.
  \item Si $X$ est compact et $f:X\to Y$ est continue et bijective, alors $f^{-1}$ est continue.
  \item Si $X$ et $Y$ sont homéomorphes et $X$ est complet, alors $Y$ est complet.
  \item Dans un espace métrique complet, une intersection dénombrable d'ouverts denses est un ouvert dense.
  \item $\Q$ est une intersection dénombrable d'ouverts de $\R$.
  \item Tout espace métrique compact est séparable.
  \item Dans un espace métrique, toute partie fermée et précompacte est compacte.
  \item Si $D$ est dense dans $X$ et $f:D\to Y$ est uniformément continue, alors $f$ se prolonge en une application continue $X\to Y$.
\end{enumerate}
\end{exo}

\begin{solution}
\noindent\textbf{1. Faux.} Sur un ensemble $X$ à au moins deux points muni de la distance discrète ($d(x,y)=1$ si $x\ne y$), $B(a,1)=\{a\}$ est fermée, donc égale à son adhérence, alors que $\Bf(a,1)=X$. On a toujours $\adh{B(a,r)}\subset\Bf(a,r)$, avec égalité dans un EVN.

\medskip
\noindent\textbf{2. Faux.} La boule unité fermée de $\ell^2(\N)$ est fermée et bornée dans un espace complet, mais la suite $(e_n)$ vérifie $\norm{e_n-e_m}=\sqrt2$ pour $n\ne m$ et n'a aucune sous-suite convergente. C'est vrai dans $\K^d$ (Corollaire~1.2), et cela caractérise la dimension finie (Théorème~3.1).

\medskip
\noindent\textbf{3. Faux.} $f:\R\to\R$ constante égale à $0$ : $f^{-1}(\{0\})=\R$. (L'image \emph{directe} d'un compact est compacte : Théorème~1.3.)

\medskip
\noindent\textbf{4. Faux.} $\arctan(\R)=\,]-\frac\pi2,\frac\pi2[$. (L'image \emph{réciproque} d'un fermé est fermée ; et l'image directe d'un fermé est fermée si l'espace de départ est compact.)

\medskip
\noindent\textbf{5. Faux.} $f(x)=1/x$ sur $]0,1]$ : la suite $(1/n)$ est de Cauchy, son image $(n)$ ne l'est pas. C'est vrai pour une application \emph{uniformément} continue, et c'est le point de départ du Théorème~1.1.

\medskip
\noindent\textbf{6. Vrai.} Soit $F$ un fermé de $X$ ; il est compact (Proposition~1.11), donc $f(F)$ est compact (Théorème~1.3), donc fermé dans $Y$ (Proposition~1.10). Or $(f^{-1})^{-1}(F)=f(F)$ : l'image réciproque par $f^{-1}$ de tout fermé est fermée, et $f^{-1}$ est continue.

\medskip
\noindent\textbf{7. Faux.} $\R$ et $]-1,1[$ sont homéomorphes par $x\mapsto\frac x{1+|x|}$, le premier est complet, le second non (Remarque~1.6). La complétude n'est pas une notion topologique.

\medskip
\noindent\textbf{8. Faux.} Dense, oui (Baire) ; ouvert, non : $\R\setminus\Q=\bigcap_{q\in\Q}(\R\setminus\{q\})$ est une intersection dénombrable d'ouverts denses de $\R$, et n'est pas ouvert puisque $\Q$ est dense.

\medskip
\noindent\textbf{9. Faux.} Si $\Q=\bigcap_nO_n$ avec $O_n$ ouverts, chaque $O_n$ contient $\Q$, donc est dense. Les ouverts $\R\setminus\{q\}$, $q\in\Q$, sont denses aussi. L'intersection de cette famille dénombrable d'ouverts denses est $\Q\cap(\R\setminus\Q)=\emptyset$, ce qui contredit le lemme de Baire dans $\R$ complet.

\medskip
\noindent\textbf{10. Vrai.} Un compact est précompact (Proposition~1.13) : pour tout $n\ge1$ il est recouvert par un nombre fini de boules de rayon $\frac1n$, de centres formant un ensemble fini $C_n$. La réunion $\bigcup_nC_n$ est dénombrable, et tout point est à distance $<\frac1n$ d'un point de $C_n$ pour tout $n$ : elle est dense.

\medskip
\noindent\textbf{11. Faux.} $X=\Q$, $A=\Q\cap[0,1]$ : fermée dans $\Q$, précompacte, mais une suite de rationnels de $[0,1]$ tendant vers $\frac1{\sqrt2}$ n'a pas de valeur d'adhérence dans $A$ (Exercice~1.21). C'est vrai si $X$ est complet (Corollaire~1.4).

\medskip
\noindent\textbf{12. Faux.} $D=Y=\Q$, $X=\R$, $f=\mathrm{id}$ : $f$ est $1$-lipschitzienne, mais un prolongement continu $\tilde f:\R\to\Q$ vérifierait $\tilde f(x)=\lim f(q_n)=x$ pour $q_n\to x$, impossible si $x\notin\Q$. Il faut $Y$ complet (Théorème~1.1).
\end{solution}

\begin{piege}
\textbf{Direct ou réciproque, topologique ou métrique.} Les questions 3, 4 et 6 : l'image \emph{réciproque} préserve ouverts et fermés, l'image \emph{directe} préserve les compacts, et rien d'autre en général. Les questions 5, 7 et 12 : la continuité est une notion topologique, la complétude et le caractère de Cauchy sont des notions métriques, et seule l'uniforme continuité fait le lien. Les questions 2 et 11 : « compact $=$ fermé borné » est propre à $\K^d$ ; l'énoncé général est « complet et précompact ». Les questions 8 et 9 : Baire donne la densité d'une intersection dénombrable d'ouverts denses, jamais son ouverture.
\end{piege}

% =====================================================================
\begin{exo}{Exercice 4 -- Vrai ou faux ? Chapitres 2 et 3}
Pour chaque affirmation, dire si elle est vraie ou fausse, et justifier par une preuve courte ou un contre-exemple. $E$ et $F$ désignent des EVN.
\begin{enumerate}[label=\arabic*.]
  \item Deux normes sur un même espace vectoriel sont toujours équivalentes.
  \item Si $E$ est de dimension finie, toute application linéaire $E\to F$ est continue.
  \item Tout sous-espace vectoriel de $E$ est fermé.
  \item Un sous-espace vectoriel strict de $E$ est d'intérieur vide.
  \item Si $T:E\to F$ est linéaire et bornée sur une boule $B(a,r)$, $r>0$, alors $T$ est continue.
  \item Pour $1\le p\le q\le\infty$, on a $L^p(]0,1[)\subset L^q(]0,1[)$.
  \item $\Cz([0,1])$ muni de la norme $\norm f_1=\int_0^1|f|$ est un espace de Banach.
  \item Un espace de Banach de dimension infinie n'admet pas de base (algébrique) dénombrable.
  \item La boule unité fermée de $\Cz([0,1])$ est compacte pour la norme uniforme.
  \item Une partie bornée et uniformément équicontinue de $\Cz([0,1])$ est compacte.
  \item Si $(f_n)$ est une suite de $\mathcal C^1([0,1])$ avec $\norm{f_n}_\infty+\norm{f_n'}_\infty\le M$ pour tout $n$, alors $(f_n)$ admet une sous-suite uniformément convergente.
  \item Si $E$ est de Banach et si $(T_n)$ est une suite de $\Lc(E,F)$ convergeant simplement vers $T$, alors $\opnorm{T_n-T}\to0$.
  \item Une limite simple d'applications linéaires continues de $E$ dans $F$ est continue.
  \item Une bijection linéaire continue $E\to F$ a un inverse continu.
\end{enumerate}
\end{exo}

\begin{solution}
\noindent\textbf{1. Faux.} Sur $\Cz([0,1])$, $\norm f_1\le\norm f_\infty$, mais $f_n(x)=x^n$ vérifie $\norm{f_n}_\infty=1$ et $\norm{f_n}_1=\frac1{n+1}\to0$ : pas de constante $C$ avec $\norm\cdot_\infty\le C\norm\cdot_1$. C'est vrai en dimension finie (Corollaire~2.1).

\medskip
\noindent\textbf{2. Vrai.} Soit $(b_1,\dots,b_d)$ une base de $E$ et $N(x):=\max_k|x_k|$ pour $x=\sum_kx_kb_k$ ; $N$ est une norme sur $E$, équivalente à $\norm\cdot_E$ (Corollaire~2.1). Alors
\[
\norm{T(x)}_F\le\sum_k|x_k|\norm{T(b_k)}_F\le\Big(\sum_k\norm{T(b_k)}_F\Big)N(x)\le C\norm x_E .
\]

\medskip
\noindent\textbf{3. Faux.} $c_{<\infty}(\N)$ est un sous-espace dense et strict de $\ell^2(\N)$ (Proposition~2.3), donc non fermé. C'est vrai pour les sous-espaces de dimension finie (Proposition~3.2).

\medskip
\noindent\textbf{4. Vrai.} Si un sous-espace $V$ contient une boule $B(a,r)$, alors pour tout $x\ne0$, $a+\frac r{2\norm x}x\in V$ et $a\in V$, donc $x\in V$ : $V=E$.

\medskip
\noindent\textbf{5. Vrai.} Si $\norm{T(x)}\le M$ sur $B(a,r)$, alors pour $\norm z<1$, $\norm{T(z)}=\frac1r\norm{T(a+rz)-T(a)}\le\frac{2M}r$ : $T$ est bornée sur la boule unité, donc continue (Théorème~2.4). C'est le calcul final de la preuve de Banach--Steinhaus.

\medskip
\noindent\textbf{6. Faux.} C'est l'inclusion inverse : sur un ensemble de mesure finie, $L^q\subset L^p$ pour $p\le q$ (Hölder). Par exemple $x^{-1/2}\in L^1(]0,1[)\setminus L^2(]0,1[)$. L'inclusion $\ell^p\subset\ell^q$ (Proposition~2.4) va dans l'autre sens, et sur $\R$ il n'y a aucune inclusion.

\medskip
\noindent\textbf{7. Faux.} Soit $f_n$ continue, affine par morceaux, nulle sur $[0,\frac12-\frac1n]$, égale à $1$ sur $[\frac12,1]$. Alors $\norm{f_n-f_m}_1\le\frac1{\min(n,m)}$ : la suite est de Cauchy pour $\norm\cdot_1$. Si elle convergeait vers $f$ continue, on aurait $\int_0^a|f|=0$ pour tout $a<\frac12$ et $\int_{1/2}^1|f-1|=0$, donc $f=0$ sur $[0,\frac12[$ et $f=1$ sur $[\frac12,1]$ : contradiction avec la continuité en $\frac12$.

\medskip
\noindent\textbf{8. Vrai.} Si $(b_n)_{n\in\N}$ était une base, $E=\bigcup_nF_n$ avec $F_n:=\Vect(b_0,\dots,b_n)$, fermés (Proposition~3.2) et d'intérieur vide (question~4, $F_n\ne E$). Cela contredit le lemme de Baire dans $E$ complet (Exercice~3.2). Ainsi $\R[x]$ n'est complet pour aucune norme.

\medskip
\noindent\textbf{9. Faux.} $\Cz([0,1])$ est de dimension infinie (Théorème~3.1). Directement : $f_n(x)=x^n$ est dans la boule unité, converge simplement vers une fonction discontinue, donc n'a pas de sous-suite uniformément convergente.

\medskip
\noindent\textbf{10. Faux.} Il manque « fermée » : l'ensemble des fonctions constantes de valeur dans $]0,1[$ est borné et uniformément équicontinu, mais non fermé, donc non compact. Une partie bornée et uniformément équicontinue est seulement d'adhérence compacte (Théorème~3.2 et Exercice~3.5).

\medskip
\noindent\textbf{11. Vrai.} La suite est bornée par $M$, et uniformément équicontinue car $|f_n(x)-f_n(y)|\le M|x-y|$ (accroissements finis). Le Corollaire~3.1 d'Arzelà--Ascoli donne une sous-suite uniformément convergente. (Sa limite est $M$-lipschitzienne, mais pas nécessairement $\mathcal C^1$.)

\medskip
\noindent\textbf{12. Faux.} Sur $E=\ell^2(\N)$, $T_n(u):=u_n$ : $T_n(u)\to0$ pour tout $u$, mais $\opnorm{T_n-0}=1$ (tester $e_n$). La convergence simple donne seulement $\sup_n\opnorm{T_n}<\infty$ et $\opnorm T\le\liminf\opnorm{T_n}$ (Corollaire~3.3).

\medskip
\noindent\textbf{13. Faux.} Sur $E=\R[x]$ muni de la norme uniforme sur $[0,1]$, les formes $A_n(P):=n\big(P(1)-P(1-\frac1n)\big)$ sont continues ($|A_n(P)|\le2n\norm P_\infty$) et convergent simplement vers $P\mapsto P'(1)$, qui est discontinue (Contre-exemple~2.1, Exercice~3.9). C'est vrai si $E$ est de Banach (Corollaire~3.3) : ici $E$ ne l'est pas (question~8).

\medskip
\noindent\textbf{14. Faux.} L'identité de $(c_{<\infty},\norm\cdot_1)$ dans $(c_{<\infty},\norm\cdot_\infty)$ est linéaire, bijective et continue ($\norm u_\infty\le\norm u_1$), mais son inverse ne l'est pas : $u^{(n)}:=e_0+\dots+e_{n-1}$ vérifie $\norm{u^{(n)}}_\infty=1$ et $\norm{u^{(n)}}_1=n$. C'est vrai entre espaces de Banach (Corollaire~3.4).
\end{solution}

\begin{piege}
\textbf{Dimension finie, complétude : les deux hypothèses cachées.} En dimension finie tout est vrai : normes équivalentes, applications linéaires continues, sous-espaces fermés, fermés bornés compacts (1, 2, 3, 9). En dimension infinie, c'est la \emph{complétude} qui remplace la dimension finie dans les grands théorèmes : limite simple d'opérateurs, inverse continu, absence de base dénombrable (8, 13, 14). Les contre-exemples à connaître par cœur sont donc des espaces non complets, $c_{<\infty}$ et $\R[x]$, ou une norme trop faible sur un bon espace, $(\Cz([0,1]),\norm\cdot_1)$ (7).
\end{piege}

\end{document}
